% Отчёт по ЛР1. Сборка: ./template/build.sh 1
\documentclass[a4paper,14pt]{extarticle}
\usepackage{../../../template/gostreport}

\begin{document}

\makegosttitle
  {1}
  {Кодирование и шифрование}
  {студент группы R3243 Осипов А.\,П.\\[4pt]
   студент группы R3238 Федорец Л.\,Б.}
  {Догадин Е.\,В.}

\tableofcontents
\newpage

\unnumsection{Введение}

В работе рассматривается применение матричного умножения к двум задачам: шифрованию
сообщений (шифр Хилла) и помехоустойчивому кодированию (код Хэмминга). Основная
особенность работы состоит в том, что все вычисления выполняются не над полем
вещественных чисел, а над кольцом вычетов $\mathbb{Z}_n$ (задания~1 и~2) и над полем
$\mathbb{F}_2$ (задание~3). Это меняет смысл привычных операций: деление заменяется
умножением на обратный элемент, а обратимость матрицы определяется не условием
$\det \ne 0$, а условием взаимной простоты определителя с модулем.

Вычисления выполнены на языке Python с использованием библиотеки \texttt{sympy}
(целочисленная матричная арифметика, вычисление определителя и присоединённой матрицы).
Обратный элемент по модулю вычисляется собственной реализацией расширенного алгоритма
Евклида. Полный текст программы приведён в приложении~А.

\section{Задание 1. Шифр Хилла}

\subsection{Алфавит и исходное сообщение}

В качестве алфавита взяты 26 букв латиницы и три служебных символа
(\texttt{\_}, \texttt{!}, \texttt{?}), всего $n = 29$ символов, пронумерованных
от~0 до~28. Число $n = 29$ простое, и это выбрано намеренно: при простом модуле
любой определитель, не кратный $n$, автоматически взаимно прост с $n$, поэтому
проверка обратимости ключа сводится к условию $\det K \not\equiv 0 \pmod{29}$.
Соответствие символов и номеров приведено в таблице~\ref{tab:alphabet}.

\begin{table}[H]
\caption{Соответствие символов алфавита и их числовых номеров}
\label{tab:alphabet}
\centering\footnotesize\setlength{\tabcolsep}{3.4pt}
\begin{tabular}{l|ccccccccccccccc}
\toprule
Символ & \texttt{A} & \texttt{B} & \texttt{C} & \texttt{D} & \texttt{E} & \texttt{F} & \texttt{G} & \texttt{H} & \texttt{I} & \texttt{J} & \texttt{K} & \texttt{L} & \texttt{M} & \texttt{N} & \texttt{O} \\
Номер & 0 & 1 & 2 & 3 & 4 & 5 & 6 & 7 & 8 & 9 & 10 & 11 & 12 & 13 & 14 \\
\midrule
Символ & \texttt{P} & \texttt{Q} & \texttt{R} & \texttt{S} & \texttt{T} & \texttt{U} & \texttt{V} & \texttt{W} & \texttt{X} & \texttt{Y} & \texttt{Z} & \texttt{\_} & \texttt{!} & \texttt{?} &  \\
Номер & 15 & 16 & 17 & 18 & 19 & 20 & 21 & 22 & 23 & 24 & 25 & 26 & 27 & 28 &  \\
\bottomrule
\end{tabular}
\end{table}

Шифруемое сообщение --- \texttt{I\_LOVE\_ITMO!} длиной 12 символов.
Длина выбрана не случайно: 12 делится на 2, 3 и~4, то есть на размеры всех трёх
ключей, и последний блок ни при одном из них не приходится дополнять служебными
символами. Числовое представление сообщения приведено в таблице~\ref{tab:message}.

\begin{table}[H]
\caption{Исходное сообщение и его числовое представление}
\label{tab:message}
\centering\footnotesize\setlength{\tabcolsep}{4.5pt}
\begin{tabular}{l|cccccccccccc}
\toprule
Символ & \texttt{I} & \texttt{\_} & \texttt{L} & \texttt{O} & \texttt{V} & \texttt{E} & \texttt{\_} & \texttt{I} & \texttt{T} & \texttt{M} & \texttt{O} & \texttt{!} \\
Номер & 8 & 26 & 11 & 14 & 21 & 4 & 26 & 8 & 19 & 12 & 14 & 27 \\
\bottomrule
\end{tabular}
\end{table}

\subsection{Метод шифрования Хилла}

Сообщение переводится в последовательность номеров символов, которая разбивается на
блоки длины $z$, где $z$ --- размерность квадратной матрицы-ключа $K$. Каждый блок
рассматривается как вектор-столбец $\vect{v} \in \mathbb{Z}_n^{z}$ и шифруется
умножением на ключ по модулю $n$:
\begin{equation}
  \vect{c} = K\vect{v} \bmod n .
  \label{eq:enc}
\end{equation}
Дешифрация выполняется умножением на обратную по модулю матрицу:
\begin{equation}
  \vect{v} = K^{-1}\vect{c} \bmod n .
  \label{eq:dec}
\end{equation}
Приведение по модулю необходимо, чтобы результат оставался в диапазоне $[0, n-1]$,
то есть соответствовал какому-либо символу алфавита.

Существенно, что блоки обрабатываются независимо друг от друга: матрица умножается
на каждый вектор отдельно, и никакого переноса информации между блоками не
происходит. Например, фрагмент \texttt{I\_LOVE} с номерами
$(8, 26, 11, 14, 21, 4)$ при ключе размера $2$ разбивается на блоки
$(8, 26)$, $(11, 14)$, $(21, 4)$, а при ключе размера $3$ --- на блоки
$(8, 26, 11)$ и $(14, 21, 4)$. Это свойство определяет и поведение шифра при
искажениях, что разбирается в подразделе~\ref{sec:corrupt}.

\subsection{Матрицы-ключи и их определители}

Выбраны три матрицы-ключа размеров $2\times2$, $3\times3$ и $4\times4$:

\begin{equation}
  K_2 = \begin{pmatrix} 1 & 1 \\ 2 & 8 \end{pmatrix}, \qquad K_3 = \begin{pmatrix} 7 & 0 & 28 \\ 4 & 4 & 4 \\ 8 & 1 & 2 \end{pmatrix},
\end{equation}
\begin{equation}
  K_4 = \begin{pmatrix} 3 & 3 & 17 & 18 \\ 22 & 9 & 1 & 15 \\ 10 & 7 & 8 & 8 \\ 20 & 11 & 7 & 20 \end{pmatrix}.
\end{equation}

Определитель матрицы $K_2$ вычисляется непосредственно:
\begin{equation}
  \det K_2 = 1\cdot 8 - 1\cdot 2 = 6 .
\end{equation}
Определители $K_3$ и $K_4$ вычислены по аналогии (разложением по строке) и равны
$\det K_3 = -756$ и $\det K_4 = 5285$. Поскольку работа ведётся в $\mathbb{Z}_{29}$,
значение имеют не сами определители, а их вычеты по модулю~29:
\begin{equation}
  \det K_2 \equiv 6, \qquad \det K_3 \equiv 27, \qquad \det K_4 \equiv 7 \pmod{29}.
\end{equation}
Все три вычета отличны от нуля и попарно различны, как того требует задание.
Так как 29 --- простое число, из $\det K \not\equiv 0$ сразу следует
$\gcd(\det K, 29) = 1$, то есть все три ключа обратимы по модулю~29.

\subsection{Шифрование сообщения}

Покажем расчёт на первом блоке для ключа $K_2$. Блок \texttt{I\_} имеет числовое
представление $(8, 26)$, и по формуле~\eqref{eq:enc}:
\begin{equation}
  K_2 \begin{pmatrix} 8 \\ 26 \end{pmatrix}
  = \begin{pmatrix} 1\cdot 8 + 1\cdot 26 \\ 2\cdot 8 + 8\cdot 26 \end{pmatrix}
  = \begin{pmatrix} 34 \\ 224 \end{pmatrix}
  \equiv \begin{pmatrix} 5 \\ 21 \end{pmatrix} \pmod{29},
\end{equation}
поскольку $34 = 1\cdot 29 + 5$ и $224 = 7\cdot 29 + 21$. Номерам 5 и 21
соответствуют символы \texttt{F} и \texttt{V}, что даёт первые два символа
шифрограммы. Остальные блоки и остальные ключи обработаны по той же схеме
программно; полученные шифрограммы приведены в таблице~\ref{tab:corrupt}.

\subsection{Обратные матрицы по модулю}

Обратная матрица по модулю $n$ вычисляется по формуле
\begin{equation}
  K^{-1} \equiv (\det K)^{-1} \cdot \operatorname{adj} K \pmod n ,
  \label{eq:inv}
\end{equation}
где $\operatorname{adj} K$ --- присоединённая матрица (транспонированная матрица
алгебраических дополнений), а $(\det K)^{-1}$ --- обратный элемент к определителю
в $\mathbb{Z}_n$. Функция \texttt{numpy.linalg.inv} здесь неприменима: она
возвращает вещественную обратную матрицу, не имеющую отношения к арифметике вычетов.

Обратный элемент находится расширенным алгоритмом Евклида. Он даёт такие целые
$x$ и $y$, что $ax + ny = \gcd(a, n)$; при $\gcd(a, n) = 1$ приведение этого
равенства по модулю $n$ обнуляет слагаемое $ny$ и оставляет $ax \equiv 1$, то есть
искомым обратным элементом является коэффициент $x$. Для $\det K_2 = 6$:
\begin{align}
  29 &= 4\cdot 6 + 5, \\
   6 &= 1\cdot 5 + 1,
\end{align}
откуда обратным ходом
\begin{equation}
  1 = 6 - 1\cdot 5 = 6 - (29 - 4\cdot 6) = 5\cdot 6 - 1\cdot 29 ,
\end{equation}
то есть $6^{-1} \equiv 5 \pmod{29}$. Аналогично получено
$(\det K_3)^{-1} \equiv 14$ и $(\det K_4)^{-1} \equiv 25$.

Присоединённая матрица для $K_2$ равна $\operatorname{adj} K_2 = \begin{pmatrix} 8 & -1 \\ -2 & 1 \end{pmatrix}$,
и по формуле~\eqref{eq:inv}
\begin{equation}
  K_2^{-1} \equiv 5 \cdot \begin{pmatrix} 8 & -1 \\ -2 & 1 \end{pmatrix} = \begin{pmatrix} 40 & -5 \\ -10 & 5 \end{pmatrix} \equiv \begin{pmatrix} 11 & 24 \\ 19 & 5 \end{pmatrix} \pmod{29}.
\end{equation}
Для остальных ключей вычисления выполнены по той же формуле программно:
\begin{equation}
  K_3^{-1} \equiv \begin{pmatrix} 27 & 15 & 27 \\ 17 & 18 & 16 \\ 14 & 18 & 15 \end{pmatrix}, \qquad K_4^{-1} \equiv \begin{pmatrix} 22 & 8 & 7 & 12 \\ 25 & 6 & 7 & 5 \\ 27 & 14 & 23 & 14 \\ 7 & 7 & 13 & 21 \end{pmatrix} \pmod{29}.
\end{equation}

Корректность обратных матриц проверена непосредственно: произведение ключа на
найденную обратную матрицу по модулю~29 должно давать единичную матрицу. Для $K_2$:
\begin{equation}
\begin{aligned}
  K_2 K_2^{-1} &= \begin{pmatrix} 1 & 1 \\ 2 & 8 \end{pmatrix}\begin{pmatrix} 11 & 24 \\ 19 & 5 \end{pmatrix} = \begin{pmatrix} 30 & 29 \\ 174 & 88 \end{pmatrix} \\
  &\equiv \begin{pmatrix} 1 & 0 \\ 0 & 1 \end{pmatrix} \pmod{29}.
\end{aligned}
\end{equation}
Для ключей $K_3$ и $K_4$ проверка выполнена программно и также даёт единичную
матрицу соответствующего размера (см.~листинг~\ref{lst:task1}).
Дешифрация шифрограмм без внесения искажений во всех трёх случаях восстанавливает
исходное сообщение \texttt{I\_LOVE\_ITMO!} посимвольно.

\subsection{Вредоносное вмешательство и дешифрация}
\label{sec:corrupt}

В каждой из трёх шифрограмм заменены символы на позициях 1, 2 и~4 (нумерация
с единицы) на символы \texttt{Z}, \texttt{X} и~\texttt{C} соответственно.
Позиции выбраны одинаковыми для всех трёх ключей намеренно: при таком выборе
единственным различием между тремя экспериментами остаётся размер блока, что
делает результаты сравнимыми между собой. В таблице~\ref{tab:corrupt} искажённые
позиции шифрограмм и отличающиеся от оригинала символы расшифровок выделены
полужирным начертанием.

\begin{table}[H]
\caption{Шифрограммы, внесённые искажения и результаты дешифрации}
\label{tab:corrupt}
\centering\footnotesize\setlength{\tabcolsep}{6pt}
\begin{tabular}{c|c|c|c}
\toprule
Ключ & Шифрограмма & После вмешательства & Результат дешифрации \\
\midrule
$K_{2}$ & \texttt{\textbf{F}\textbf{V}Z\textbf{S}ZQFACSMM} & \texttt{\textbf{Z}\textbf{X}Z\textbf{C}ZQFACSMM} & \texttt{\textbf{P}\textbf{K}\textbf{E}\textbf{V}VE\_ITMO!} \\
$K_{3}$ & \texttt{\textbf{Q}\textbf{G}Z\textbf{H}LZSJW?JT} & \texttt{\textbf{Z}\textbf{X}Z\textbf{C}LZSJW?JT} & \texttt{\textbf{N}\textbf{V}\textbf{I}\textbf{Y}\textbf{X}\textbf{V}\_ITMO!} \\
$K_{4}$ & \texttt{\textbf{T}\textbf{W}!\textbf{U}XGRXFRWT} & \texttt{\textbf{Z}\textbf{X}!\textbf{C}XGRXFRWT} & \texttt{\textbf{T}\textbf{F}\textbf{W}\textbf{E}VE\_ITMO!} \\
\bottomrule
\end{tabular}
\end{table}

Сводка по числу пострадавших символов приведена в таблице~\ref{tab:damage}.

\begin{table}[H]
\caption{Зависимость числа искажённых символов от размера блока}
\label{tab:damage}
\centering\footnotesize\setlength{\tabcolsep}{6pt}
\begin{tabular}{c|c|c|c|c}
\toprule
Ключ & Размер блока & Всего блоков & Задето блоков & Искажено символов \\
\midrule
$K_{2}$ & 2 & 6 & 2 & 4 \\
$K_{3}$ & 3 & 4 & 2 & 6 \\
$K_{4}$ & 4 & 3 & 1 & 4 \\
\bottomrule
\end{tabular}
\end{table}

\subsection*{Вывод по заданию~1}
\addcontentsline{toc}{subsection}{Вывод по заданию~1}

Шифрование и дешифрация выполнены для всех трёх ключей. При отсутствии вмешательства
обратное преобразование восстанавливает исходное сообщение полностью, то есть метод
Хилла работает корректно, а построенные обратные по модулю матрицы найдены верно ---
это же подтверждается равенством $K K^{-1} \equiv E$.

Основной результат касается распространения ошибок. Искажение шифрограммы не выходит
за пределы того блока, в который попал испорченный символ: остальная часть сообщения
расшифровывается без потерь. Причина в том, что при шифровании каждый блок умножается
на матрицу независимо от других, и информация между блоками не смешивается.

При этом число пострадавших символов определяется не количеством ошибок, а тем, как
они распределились по блокам. Три искажения в одних и тех же позициях дали разный
результат: 4 испорченных символа для ключа $K_2$ и столько же для $K_4$, но 6 символов
для $K_3$. Для $K_2$ ошибки задели два блока по два символа, для $K_3$ --- тоже два
блока, но уже по три символа, а для $K_4$ все три ошибки пришлись на один
четырёхсимвольный блок. Отсюда следует, что с ростом размера блока растёт и цена
одиночной ошибки, зато уменьшается вероятность задеть несколько блоков сразу.

Внутри задетого блока портятся, вообще говоря, все символы. Позиция $i$ расшифрованного
блока искажается тогда, когда элемент $(K^{-1})_{ij}$ отличен от нуля, где $j$ ---
позиция испорченного символа. Отсюда видно, при каком условии одна ошибка приводила бы
ровно к одной ошибке в расшифровке: в каждой строке обратной матрицы должен стоять
ровно один ненулевой элемент. Однако такой ключ обесценивает шифр --- каждый символ
всего лишь умножается на свою константу и, возможно, меняет позицию, смешивания между
символами не происходит. Шифр вырождается в моноалфавитную подстановку, которая
вскрывается частотным анализом.

\section{Задание 2. Взлом шифра Хилла}

\subsection{Постановка}

Моделируется следующая ситуация: имеются две шифрограммы, полученные шифром Хилла
с одним и тем же ключом, и этот ключ неизвестен. Однако в руки попала расшифровка
(оригинал) одной из них. Требуется найти способ прочитать вторую.

\subsection{Исходные данные}

Используется алфавит из задания~1 ($n = 29$). Ключ размера $2\times2$ генерируется
программой случайно: каждый элемент берётся из диапазона $[0, 28]$, после чего
проверяется условие $\det K \not\equiv 0 \pmod{29}$, и при его нарушении матрица
генерируется заново. Значение ключа при выводе не используется --- он известен
только программе. Для воспроизводимости расчётов зафиксировано начальное значение
генератора: \texttt{random.seed(67)}.

Взяты два сообщения по 12 символов, \texttt{PRETTY\_GIRLS} и
\texttt{STAY\_WITH\_ME}, с числовыми представлениями

\begin{equation}
\begin{aligned}
  \vect{m}_1 &= (15,\ 17,\ 4,\ 19,\ 19,\ 24,\ 26,\ 6,\ 8,\ 17,\ 11,\ 18), \\
  \vect{m}_2 &= (18,\ 19,\ 0,\ 24,\ 26,\ 22,\ 8,\ 19,\ 7,\ 26,\ 12,\ 4).
\end{aligned}
\end{equation}
Шифрование обоих сообщений неизвестным ключом даёт шифрограммы
\texttt{XBHYXOMGJHSF} и \texttt{GMOWCLPEFNHV}.

После этого первое сообщение считается забытым (в программе оно заменено на
строку из двенадцати символов \texttt{?}). Атакующему доступны: обе шифрограммы
и открытый текст второго сообщения. Ключ недоступен.

\subsection{Вывод способа нахождения ключа}

Дешифрация одного блока записывается как
\begin{equation}
  D\vect{x} \equiv \vect{b} \pmod n ,
  \label{eq:dec2}
\end{equation}
где $D$ --- матрица дешифрования, $\vect{x}$ --- блок шифрограммы,
$\vect{b}$ --- соответствующий блок открытого текста. Неизвестной здесь
является именно $D$, поэтому её оставляем слева, а $\vect{x}$ переносим вправо,
где она обращается.

Одного блока для этого недостаточно. Умножение блока на матрицу $2\times2$ даёт
два уравнения, тогда как в ключе четыре неизвестных элемента --- система
недоопределена. Поэтому вместо вектора $\vect{x}$ берётся матрица $X$, столбцами
которой служат два блока шифрограммы, и точно так же из соответствующих блоков
открытого текста составляется матрица $B$. Два блока дают четыре уравнения,
чего ровно хватает. Выбранные блоки при этом не должны быть линейно зависимы
по модулю $n$, иначе матрица $X$ необратима и решение не существует.
Уравнение~\eqref{eq:dec2} превращается в матричное
\begin{equation}
  D X \equiv B \pmod n \qquad\Longrightarrow\qquad D \equiv B X^{-1} \pmod n .
  \label{eq:attack}
\end{equation}
Исходный ключ получается ещё одним обращением: $K = D^{-1} \bmod n$.

\subsection{Нахождение ключа}

Берутся первые два блока. Из шифрограммы второго сообщения составляется
\begin{equation}
  X = \begin{pmatrix} 6 & 14 \\ 12 & 22 \end{pmatrix}, \qquad \det X = -36 \equiv 22 \pmod{29},
\end{equation}
определитель отличен от нуля, значит блоки линейно независимы и матрица обратима.
Обратный элемент к определителю равен $22^{-1} \equiv 4$, откуда
\begin{equation}
  X^{-1} \equiv 4 \cdot \begin{pmatrix} 22 & -14 \\ -12 & 6 \end{pmatrix} \equiv \begin{pmatrix} 1 & 2 \\ 10 & 24 \end{pmatrix} \pmod{29}.
\end{equation}
Из открытого текста второго сообщения составляется матрица тех же двух блоков
$B = \begin{pmatrix} 18 & 0 \\ 19 & 24 \end{pmatrix}$, и по формуле~\eqref{eq:attack}
\begin{equation}
  D \equiv \begin{pmatrix} 18 & 0 \\ 19 & 24 \end{pmatrix}\begin{pmatrix} 1 & 2 \\ 10 & 24 \end{pmatrix} \equiv \begin{pmatrix} 18 & 7 \\ 27 & 5 \end{pmatrix} \pmod{29}.
\end{equation}
Обращение найденной матрицы даёт сам ключ:
\begin{equation}
  K = D^{-1} \equiv \begin{pmatrix} 2 & 3 \\ 24 & 13 \end{pmatrix} \pmod{29}.
\end{equation}

\subsection{Проверка и расшифровка первого сообщения}

Проверка выполняется независимо от того, как ключ был получен: известный открытый
текст второго сообщения шифруется \emph{найденным} ключом $K$, и результат
сравнивается с имеющейся шифрограммой. Оба значения совпадают:

\begin{equation}
  \text{\texttt{GMOWCLPEFNHV}} \quad = \quad \text{\texttt{GMOWCLPEFNHV}} .
\end{equation}
Совпадение не сводится к построению: в выводе ключа участвовали только первые
два блока, а сходятся все шесть.

Расшифровка первой шифрограммы найденной матрицей $D$ даёт
\texttt{PRETTY\_GIRLS} --- забытое сообщение восстановлено.

\subsection*{Вывод по заданию~2}
\addcontentsline{toc}{subsection}{Вывод по заданию~2}

Шифр Хилла удалось взломать, имея на руках одну шифрограмму вместе с её открытым
текстом. Причина в линейности преобразования: если два сообщения зашифрованы одним
ключом и рассекречены обе версии одного из них --- открытая и зашифрованная, ---
то пары «блок открытого текста~--- блок шифрограммы» дают систему линейных
уравнений относительно элементов ключа, а обратимость матриц позволяет выразить
неизвестную матрицу в явном виде.

Для ключа размера $z \times z$ требуется $z^2$ символов открытого текста: в ключе
$z^2$ неизвестных, а каждый блок даёт лишь $z$ уравнений. При меньшем объёме
известного открытого текста система остаётся недоопределённой и ключ однозначно
не восстанавливается.

Отсюда следует, что шифр проваливается при утечке всего $z^2$ расшифрованных
символов: они позволяют раскрыть ключ целиком, а вместе с ним и всю переписку,
зашифрованную тем же ключом. По меркам современных шифров это недопустимо ---
там знание пар «открытый текст~--- шифрограмма» ключа не выдаёт. Устранить
уязвимость можно, отказавшись от постоянного ключа и используя для каждого
сообщения свой.

\section{Задание 3. Код Хэмминга}

В этом задании решается задача, обратная по смыслу заданиям~1 и~2: не скрыть
содержание сообщения, а защитить его от искажений при передаче. Меняются и исходные
данные, и арифметика. Алфавит берётся другой --- русский, из 32 букв, --- и вычисления
ведутся не в кольце вычетов $\mathbb{Z}_{29}$, как в первых двух заданиях, а в поле
$\mathbb{F}_2$, где сложение совпадает с операцией <<исключающее или>>, то есть
$1 + 1 = 0$. Ни алфавит, ни модуль из заданий~1 и~2 далее не используются.

\subsection{Алфавит и исходное сообщение}

Взят русский алфавит без буквы <<Ё>>, всего 32 буквы. Число $32 = 2^5$, поэтому
каждой букве соответствует ровно пятибитовый двоичный номер и ни одна комбинация из
пяти бит не пропадает. Буква <<Ё>> исключена намеренно: с ней алфавит содержал бы 33
буквы, что потребовало бы уже шести бит на букву, и 31 комбинация из 64 осталась бы
неиспользованной. Соответствие букв, номеров и двоичных кодов приведено
в таблице~\ref{tab:alphabet3}.

\begin{table}[H]
\caption{Соответствие букв алфавита, их номеров и пятибитовых кодов}
\label{tab:alphabet3}
\centering\footnotesize\setlength{\tabcolsep}{4pt}
\begin{tabular}{l|cccccccc}
\toprule
Буква & А & Б & В & Г & Д & Е & Ж & З \\
Номер & 0 & 1 & 2 & 3 & 4 & 5 & 6 & 7 \\
Код & \texttt{00000} & \texttt{00001} & \texttt{00010} & \texttt{00011} & \texttt{00100} & \texttt{00101} & \texttt{00110} & \texttt{00111} \\
\midrule
Буква & И & Й & К & Л & М & Н & О & П \\
Номер & 8 & 9 & 10 & 11 & 12 & 13 & 14 & 15 \\
Код & \texttt{01000} & \texttt{01001} & \texttt{01010} & \texttt{01011} & \texttt{01100} & \texttt{01101} & \texttt{01110} & \texttt{01111} \\
\midrule
Буква & Р & С & Т & У & Ф & Х & Ц & Ч \\
Номер & 16 & 17 & 18 & 19 & 20 & 21 & 22 & 23 \\
Код & \texttt{10000} & \texttt{10001} & \texttt{10010} & \texttt{10011} & \texttt{10100} & \texttt{10101} & \texttt{10110} & \texttt{10111} \\
\midrule
Буква & Ш & Щ & Ъ & Ы & Ь & Э & Ю & Я \\
Номер & 24 & 25 & 26 & 27 & 28 & 29 & 30 & 31 \\
Код & \texttt{11000} & \texttt{11001} & \texttt{11010} & \texttt{11011} & \texttt{11100} & \texttt{11101} & \texttt{11110} & \texttt{11111} \\
\bottomrule
\end{tabular}
\end{table}

Кодируемое слово --- <<СОНЯ>>. Его побуквенное двоичное представление приведено
в таблице~\ref{tab:message3}.

\begin{table}[H]
\caption{Исходное слово и его двоичное представление}
\label{tab:message3}
\centering\footnotesize\setlength{\tabcolsep}{8pt}
\begin{tabular}{l|cccc}
\toprule
Буква & С & О & Н & Я \\
Номер & 17 & 14 & 13 & 31 \\
Код & \texttt{10001} & \texttt{01110} & \texttt{01101} & \texttt{11111} \\
\bottomrule
\end{tabular}
\end{table}

Склеивание кодов даёт двоичную последовательность из 20 символов:
\begin{equation}
  \texttt{10001011100110111111}.
  \label{eq:msg20}
\end{equation}

Число 20 делится на~4 нацело, поэтому при разбиении на блоки по четыре бита получается
ровно пять блоков и последний не приходится дополнять нулями. В отличие от задания~1,
здесь это не результат выбора: длина слова (четыре буквы) задана условием, ширина
буквы (пять бит) определяется размером алфавита, а произведение $4 \cdot 5 = 20$
просто оказалось кратно четырём. При слове из трёх букв получилось бы 15 бит, и
последний блок пришлось бы дополнять. Существенное же отличие от задания~1 в другом:
там размер блока выбирался вместе с размером ключа, здесь он жёстко задан самим
кодом --- код Хэмминга $(7,4)$ принимает на вход ровно четыре бита.

\subsection{Построение матриц G и H}

Если дописать к сообщению один бит чётности, можно установить факт наличия ошибки,
но не её положение. Чтобы определить положение, проверок должно быть несколько, и их
совокупный результат должен различать позиции. Три проверки дают три бита ответа,
то есть 8 различных значений. Одно значение ($000$) отводится под случай <<ошибок
нет>>, остаются 7~--- ровно столько позиций и имеет смысл брать в кодовом слове.
Отсюда длина 7: большее число позиций три проверки уже не различат, меньшее оставит
часть значений неиспользованной. Из семи позиций три занимают сами проверочные биты,
четыре остаются информационными, что и даёт обозначение $(7,4)$.

Позиции кодового слова нумеруются числами от~1 до~7. Под проверочные биты отводятся
позиции, являющиеся степенями двойки: 1, 2 и~4; остальные позиции --- 3, 5, 6, 7 ---
занимают информационные биты $d_1, d_2, d_3, d_4$. Проверочные биты нумеруются не
подряд, а по номеру занимаемой позиции: $p_1$, $p_2$, $p_4$. При таком обозначении
всё правило формулируется одной фразой: проверочный бит $p_k$ стоит на позиции $k$ и
отвечает за бит веса $k$ в двоичной записи номера позиции. Кодовое слово имеет вид
\begin{equation}
  \vect{c} = (p_1,\ p_2,\ d_1,\ p_4,\ d_2,\ d_3,\ d_4)^{\mathsf{T}} .
  \label{eq:cword}
\end{equation}

Проверка $p_k$ охватывает те позиции, в двоичной записи номера которых присутствует
бит веса $k$. Это даёт три группы и три формулы:
\begin{equation}
\begin{aligned}
  p_1 &\ \text{--- позиции } 1,3,5,7: & p_1 &= d_1 + d_2 + d_4, \\
  p_2 &\ \text{--- позиции } 2,3,6,7: & p_2 &= d_1 + d_3 + d_4, \\
  p_4 &\ \text{--- позиции } 4,5,6,7: & p_4 &= d_2 + d_3 + d_4 .
\end{aligned}
\label{eq:parity}
\end{equation}
Все суммы берутся по модулю~2. Каждая проверка настроена так, чтобы сумма всех
охваченных ею битов, включая сам проверочный, равнялась нулю.

Выбор степеней двойки под проверочные биты объясняется именно этим: позиция~1 входит
только в первую группу, позиция~2 --- только во вторую, позиция~4 --- только в третью.
Поэтому каждый проверочный бит вычисляется непосредственно из информационных и не
входит в чужие уравнения. При размещении проверочных битов на любых других позициях
они оказались бы связаны друг с другом, и систему пришлось бы решать отдельно.

Порождающая матрица $G$ представляет собой запись формул~\eqref{eq:parity} и
переноса информационных битов по строкам; столбцов в ней четыре --- по числу
информационных битов. Кодирование выполняется умножением
\begin{equation}
  \vect{c} = G\vect{d} \bmod 2 ,
  \label{eq:henc}
\end{equation}
\begin{equation}
  G = \begin{pmatrix} 1&1&0&1 \\ 1&0&1&1 \\ 1&0&0&0 \\ 0&1&1&1 \\ 0&1&0&0 \\ 0&0&1&0 \\ 0&0&0&1 \end{pmatrix}.
  \label{eq:G}
\end{equation}
Строки сверху вниз соответствуют позициям $p_1$, $p_2$, $d_1$, $p_4$, $d_2$, $d_3$,
$d_4$. Строки 3, 5, 6 и~7 --- это строки единичной матрицы, они переносят
информационные биты без изменений; строки 1, 2 и~4 содержат коэффициенты трёх
проверок.

Проверочная матрица $H$ строится ещё проще: в столбец с номером $j$ записывается
двоичное представление числа $j$, старший бит сверху:
\begin{equation}
  H = \begin{pmatrix} 0&0&0&1&1&1&1 \\ 0&1&1&0&0&1&1 \\ 1&0&1&0&1&0&1 \end{pmatrix}.
  \label{eq:H}
\end{equation}
Её столбцы --- это числа $001, 010, 011, 100, 101, 110, 111$, то есть все числа от~1
до~7, и каждый ненулевой вектор пространства $\mathbb{F}_2^3$ встречается ровно один
раз. Строки при этом автоматически оказываются теми самыми проверками: в нижней строке
единицы стоят на позициях 1, 3, 5, 7 (номера, содержащие бит веса~1), в средней ---
на позициях 2, 3, 6, 7, в верхней --- на позициях 4, 5, 6, 7. Таким образом, $H$ есть
матричная запись трёх уравнений проверки, а условие $H\vect{c} = \vect{0}$ означает,
что все три проверки сошлись.

Порядок строк $H$ выбран так, что сверху стоит бит веса~4, а снизу --- бит веса~1.
Благодаря этому синдром, прочитанный сверху вниз как двоичное число, сразу даёт номер
испорченной позиции, и никакой таблицы соответствия не требуется.

Согласованность матриц проверяется равенством
\begin{equation}
  HG \equiv 0 \pmod 2 ,
  \label{eq:HG}
\end{equation}
где справа стоит нулевая матрица размера $3 \times 4$; вычисление выполнено программно
(см.~листинг~\ref{lst:task3}). Это равенство и означает, что любое кодовое слово
$\vect{c} = G\vect{d}$ проходит все три проверки.

\subsection{Кодирование сообщения}

Последовательность~\eqref{eq:msg20} разбивается на пять блоков по четыре бита:
\texttt{1000}, \texttt{1011}, \texttt{1001}, \texttt{1011}, \texttt{1111}. Каждый блок
рассматривается как вектор-столбец $\vect{d}$ и умножается на матрицу $G$ по
формуле~\eqref{eq:henc}. Покажем расчёт на третьем блоке $\vect{d} = (1, 0, 0, 1)$:
\begin{equation}
  G\vect{d} =
  \begin{pmatrix} 1&1&0&1 \\ 1&0&1&1 \\ 1&0&0&0 \\ 0&1&1&1 \\ 0&1&0&0 \\ 0&0&1&0 \\ 0&0&0&1 \end{pmatrix}
  \begin{pmatrix} 1 \\ 0 \\ 0 \\ 1 \end{pmatrix}
  =
  \begin{pmatrix} 1+0+0+1 \\ 1+0+0+1 \\ 1+0+0+0 \\ 0+0+0+1 \\ 0+0+0+0 \\ 0+0+0+0 \\ 0+0+0+1 \end{pmatrix}
  =
  \begin{pmatrix} 2 \\ 2 \\ 1 \\ 1 \\ 0 \\ 0 \\ 1 \end{pmatrix}
  \equiv
  \begin{pmatrix} 0 \\ 0 \\ 1 \\ 1 \\ 0 \\ 0 \\ 1 \end{pmatrix} \pmod 2 ,
\end{equation}
то есть кодовое слово третьего блока равно \texttt{0011001}. Проверочные биты при этом
получаются равными $p_1 = d_1 + d_2 + d_4 = 1 + 0 + 1 \equiv 0$,
$p_2 = d_1 + d_3 + d_4 = 1 + 0 + 1 \equiv 0$ и
$p_4 = d_2 + d_3 + d_4 = 0 + 0 + 1 = 1$, что согласуется с
формулами~\eqref{eq:parity}. Остальные блоки обработаны по той же схеме программно;
результаты приведены в таблице~\ref{tab:blocks3}.

\begin{table}[H]
\caption{Кодирование блоков сообщения матрицей $G$}
\label{tab:blocks3}
\centering\footnotesize\setlength{\tabcolsep}{8pt}
\begin{tabular}{c|c|c}
\toprule
Блок & Информационные биты & Кодовое слово \\
\midrule
1 & \texttt{1000} & \texttt{1110000} \\
2 & \texttt{1011} & \texttt{0110011} \\
3 & \texttt{1001} & \texttt{0011001} \\
4 & \texttt{1011} & \texttt{0110011} \\
5 & \texttt{1111} & \texttt{1111111} \\
\bottomrule
\end{tabular}
\end{table}

Закодированное сообщение длиной 35 бит:
\begin{equation}
  \texttt{1110000\,0110011\,0011001\,0110011\,1111111}.
  \label{eq:coded35}
\end{equation}
Длина выросла с 20 до 35 бит, то есть в $7/4$ раза. Это и есть плата за возможность
исправлять ошибки: три дополнительных бита на каждые четыре информационных.

\subsection{Вредоносное вмешательство}

Биты закодированного сообщения нумеруются с единицы сквозной нумерацией по всем 35
позициям; блок с номером $b$ занимает позиции с $7(b-1)+1$ по $7b$.

Проведены две серии по четыре опыта. В каждой серии последовательно искажаются 1, 2,
3 и~4 бита, причём каждый следующий опыт добавляет ошибку к предыдущим. Серии
отличаются только распределением ошибок по блокам, и именно это различие является
предметом исследования.

\textbf{Серия A.} Все ошибки накапливаются внутри одного блока --- третьего,
занимающего позиции 15--21. Искажаются позиции 17, затем 17 и~19, затем 17, 19 и~20,
затем 17, 19, 20 и~21; внутри блока это локальные позиции 3, 5, 6 и~7. Закодированное
сообщение с выделенными искажаемыми битами:
\begin{equation}
  \texttt{1110000\,0110011\,00\textbf{1}1\textbf{001}\,0110011\,1111111}.
\end{equation}
Соответствующие изменения третьего блока приведены в таблице~\ref{tab:spoilA}.

\begin{table}[H]
\caption{Серия A: искажение битов внутри одного блока}
\label{tab:spoilA}
\centering\footnotesize\setlength{\tabcolsep}{8pt}
\begin{tabular}{c|c|c|c}
\toprule
Опыт & Глобальные позиции & Локальные позиции & Третий блок \\
\midrule
1 бит & 17 & 3 & \texttt{0011001} $\to$ \texttt{0001001} \\
2 бита & 17, 19 & 3, 5 & \texttt{0011001} $\to$ \texttt{0001101} \\
3 бита & 17, 19, 20 & 3, 5, 6 & \texttt{0011001} $\to$ \texttt{0001111} \\
4 бита & 17, 19, 20, 21 & 3, 5, 6, 7 & \texttt{0011001} $\to$ \texttt{0001110} \\
\bottomrule
\end{tabular}
\end{table}

\textbf{Серия B.} Столько же ошибок, но каждая очередная вносится в новый блок.
Искажаются позиции 3, 12, 22 и~35:
\begin{equation}
  \texttt{11\textbf{1}0000\,0110\textbf{0}11\,0011001\,\textbf{0}110011\,111111\textbf{1}}.
\end{equation}
Распределение ошибок по блокам приведено в таблице~\ref{tab:spoilB}. Третья ошибка
намеренно внесена не в информационный, а в проверочный бит --- чтобы проверить, как
код ведёт себя при повреждении собственной служебной части.

\begin{table}[H]
\caption{Серия B: распределение ошибок по разным блокам}
\label{tab:spoilB}
\centering\footnotesize\setlength{\tabcolsep}{6pt}
\begin{tabular}{c|c|c|c|c|c}
\toprule
Опыт & Позиция & Блок & Локальная позиция & Бит & Задето блоков \\
\midrule
1 бит & 3 & 1 & 3 & $d_1$, информационный & 1 \\
2 бита & 12 & 2 & 5 & $d_2$, информационный & 2 \\
3 бита & 22 & 4 & 1 & $p_1$, проверочный & 3 \\
4 бита & 35 & 5 & 7 & $d_4$, информационный & 4 \\
\bottomrule
\end{tabular}
\end{table}

\subsection{Декодирование по синдрому}

Принятый блок представляется как $\vect{r} = \vect{c} + \vect{e}$, где $\vect{c}$ ---
переданное кодовое слово, а $\vect{e}$ --- вектор ошибок, содержащий единицы на тех
позициях, где биты изменились. Синдром вычисляется как
\begin{equation}
  \vect{s} = H\vect{r} = H(\vect{c} + \vect{e}) = H\vect{c} + H\vect{e} = H\vect{e}
  \pmod 2 ,
  \label{eq:syn}
\end{equation}
поскольку $H\vect{c} = \vect{0}$ по построению. В этом и состоит основное свойство
метода: синдром не зависит от передаваемого сообщения и определяется исключительно
вектором ошибок. Декодеру не требуется знать, что именно было отправлено.

Если ошибка ровно одна и находится на позиции $j$, то $\vect{e}$ --- столбец единичной
матрицы, и произведение $H\vect{e}$ равно $j$-му столбцу матрицы $H$, который по
построению совпадает с двоичной записью числа $j$. Прочитанный сверху вниз синдром
непосредственно даёт номер испорченного бита, а значение $\vect{s} = (0,0,0)$ означает,
что все проверки сошлись и исправлять нечего.

После исправления информационные биты извлекаются с позиций 3, 5, 6 и~7. В программе
это выполнено умножением на матрицу выборки
\begin{equation}
  R = \begin{pmatrix} 0&0&1&0&0&0&0 \\ 0&0&0&0&1&0&0 \\ 0&0&0&0&0&1&0 \\ 0&0&0&0&0&0&1 \end{pmatrix},
  \qquad RG \equiv E \pmod 2 .
  \label{eq:R}
\end{equation}
Равенство $RG = E$ проверено программно; оно означает, что из неискажённого кодового
слова исходный блок извлекается в точности в том виде, в каком он в него вошёл.
Полученные 20 бит разбиваются по пять и переводятся в буквы по
таблице~\ref{tab:alphabet3}.

Результаты серии A приведены в таблице~\ref{tab:resA}.

\begin{table}[H]
\caption{Серия A: синдромы, исправленный код и результат декодирования}
\label{tab:resA}
\centering\footnotesize\setlength{\tabcolsep}{7pt}
\begin{tabular}{c|c|c|c|c}
\toprule
Опыт & Синдром & Указанный бит & Третий блок после исправления & Слово \\
\midrule
1 бит & \texttt{011} & 3 & \texttt{0011001} & СОНЯ \\
2 бита & \texttt{110} & 6 & \texttt{0001111} & СНЭЯ \\
3 бита & \texttt{000} & ошибок нет & \texttt{0001111} & СНЭЯ \\
4 бита & \texttt{111} & 7 & \texttt{0001111} & СНЭЯ \\
\bottomrule
\end{tabular}
\end{table}

Декодированные последовательности: в первом опыте
\texttt{10001011100110111111}, то есть исходное сообщение~\eqref{eq:msg20};
в остальных трёх --- \texttt{10001011011110111111}. У четырёх остальных блоков синдром
во всех опытах серии A равен $000$, после исправления они остаются неизменными
(\texttt{1110000}, \texttt{0110011}, \texttt{0110011}, \texttt{1111111}) и
декодируются без потерь. В колонке исправленного кода видно то, что подробно
разбирается ниже: три неудачных опыта приходят в одну и ту же точку \texttt{0001111}.

Результаты серии B приведены в таблице~\ref{tab:resB}.

\begin{table}[H]
\caption{Серия B: синдромы по блокам и результат декодирования}
\label{tab:resB}
\centering\footnotesize\setlength{\tabcolsep}{7pt}
\begin{tabular}{c|c|c|c}
\toprule
Опыт & Синдромы блоков 1--5 & Код после исправления & Слово \\
\midrule
1 бит & \texttt{011, 000, 000, 000, 000} & совпадает с~\eqref{eq:coded35} & СОНЯ \\
2 бита & \texttt{011, 101, 000, 000, 000} & совпадает с~\eqref{eq:coded35} & СОНЯ \\
3 бита & \texttt{011, 101, 000, 001, 000} & совпадает с~\eqref{eq:coded35} & СОНЯ \\
4 бита & \texttt{011, 101, 000, 001, 111} & совпадает с~\eqref{eq:coded35} & СОНЯ \\
\bottomrule
\end{tabular}
\end{table}

Совпадение с~\eqref{eq:coded35} проверено побитово: восстанавливаются все пять блоков
целиком, а не только правильные буквы на выходе. Таким образом, в серии B все четыре
опыта, включая опыт с четырьмя искажёнными битами, восстановили исходное слово
полностью. Во всех опытах декодированная последовательность равна~\eqref{eq:msg20}.

Отдельного разбора заслуживает третий опыт серии B, в котором искажён проверочный бит
$p_1$ четвёртого блока. Синдром оказался равен $001 = 1$, декодер указал на первый бит
и исправил его точно так же, как исправлял информационные. Это отвечает на вопрос,
зачем проверять сами проверочные биты: декодер заранее не знает, куда попала ошибка,
для него все семь позиций равноправны, и матрица $H$ не делает между ними никакого
различия --- в её столбцах стоят просто номера от~1 до~7. Одновременно видно, что
ошибка в проверочном бите не угрожает данным: даже если оставить её на месте,
информационные биты извлекаются с позиций 3, 5, 6, 7 и не пострадают. Однако
установить, что ошибка находится именно в проверочном бите, можно лишь после вычисления
синдрома, поэтому исправление выполняется в любом случае.

\subsection{Разбор результатов серии A}

\textbf{Одна ошибка} (позиция 17, локальная~3). Синдром равен $011 = 3$, что в точности
совпадает с номером искажённого бита. Декодер переворачивает его обратно и получает
исходное кодовое слово; слово <<СОНЯ>> восстановлено полностью.

\textbf{Две ошибки} (локальные позиции 3 и~5). Синдром равен сумме двух столбцов
матрицы $H$:
\begin{equation}
  \vect{s} = \vect{h}_3 + \vect{h}_5 = 011 + 101 = 110 = 6 .
\end{equation}
Декодер трактует это как одиночную ошибку в шестом бите и переворачивает шестой бит.
В результате в блоке оказывается три искажённых бита --- 3, 5 и~6 --- вместо двух:
декодер не только не исправил ошибку, но и добавил третью. Существенно здесь не то,
что не было подано сигнала тревоги, а то, чем декодер это обосновал: на выходе у него
блок \texttt{0001111}, а это законное кодовое слово, поскольку
$\texttt{0001111} = G\,(0,1,1,1)^{\mathsf{T}}$, синдром у него нулевой, и любая
повторная проверка его пропустит. Отличить такой результат от истинного невозможно не
из-за несовершенства декодера, а потому, что необходимой для этого информации в
принятом блоке уже не осталось: получилось слово, которое имело полное право быть
переданным.

\textbf{Три ошибки} (локальные позиции 3, 5, 6). Синдром равен $000$, декодер сообщает
об отсутствии ошибок и выдаёт данные как есть. Это наихудший случай --- ошибка остаётся
незамеченной. Причина в том, что сам набор позиций $\{3, 5, 6\}$ образует кодовое
слово: $\vect{e} = \texttt{0010110} = G\,(1,1,1,0)^{\mathsf{T}}$. Сумма двух кодовых
слов снова является кодовым словом, а на кодовых словах синдром равен нулю по
определению. Три ошибки перевели принятый вектор точно в соседнее законное слово.

\textbf{Четыре ошибки} (локальные позиции 3, 5, 6, 7). Синдром равен
$\vect{h}_3 + \vect{h}_5 + \vect{h}_6 + \vect{h}_7 = 111 = 7$, декодер переворачивает
седьмой бит и тем самым устраняет четвёртую ошибку, оставаясь при том же наборе
$\{3, 5, 6\}$, что и в предыдущем опыте. Поэтому результат побитово совпадает с
трёхошибочным.

Во всех трёх неудачных опытах серии A получено одно и то же слово <<СНЭЯ>>, и это не
случайность: как видно из таблицы~\ref{tab:resA}, все три случая сходятся в одно
кодовое слово \texttt{0001111}. Декодер по построению всегда выдаёт какое-либо кодовое
слово --- других вариантов у него нет, --- и при числе ошибок больше одной он выдаёт
не переданное слово, а ближайшее к принятому. В опытах 2, 3 и~4 ближайшим каждый раз
оказалось \texttt{0001111}, откуда и одинаковый ответ.

Отдельно следует объяснить, почему пострадали именно буквы <<О>> и <<Н>>. Третий блок
кодирует биты 9--12 двадцатибитной последовательности, тогда как вторая буква занимает
биты 6--10, а третья --- биты 11--15. Четырёхбитный блок лёг на стык двух букв: биты 9
и~10 относятся ко второй букве, биты 11 и~12 --- к третьей. Поэтому искажение одного
блока из четырёх бит испортило сразу две буквы из четырёх: О~$\to$~Н и Н~$\to$~Э.

Это частный случай общего соотношения. Блок кода составляет 4 бита, буква --- 5 бит,
а числа 4 и~5 взаимно просты, поэтому их границы совпадают только через
$\operatorname{lcm}(4,5) = 20$ бит, то есть ровно один раз на всё сообщение. Из пяти
блоков целиком внутри одной буквы лежат только первый (биты 1--4 при букве, занимающей
биты 1--5) и пятый (биты 17--20 при букве, занимающей биты 16--20); три средних блока
расположены на стыке и захватывают по две буквы каждый.

Однако покрытие задаёт лишь верхнюю границу ущерба, а не сам ущерб: число реально
испорченных букв зависит ещё и от того, какие биты внутри блока изменились после
ложного исправления. В рассмотренных опытах ошибки в локальных позициях 3 и~5 после
исправления дают набор $\{3, 5, 6\}$, то есть переворачивают биты $d_1$, $d_2$, $d_3$
и не затрагивают $d_4$: последний бит блока каждый раз остаётся неизменным. Результаты
применения того же набора ошибок к каждому из пяти блоков приведены
в таблице~\ref{tab:letters}.

\begin{table}[H]
\caption{Число испорченных букв при одном и том же наборе ошибок в разных блоках}
\label{tab:letters}
\centering\footnotesize\setlength{\tabcolsep}{8pt}
\begin{tabular}{c|c|c|c}
\toprule
Испорченный блок & Покрываемые буквы & Испорчено букв & Слово \\
\midrule
1 & 1 & 1 & НОНЯ \\
2 & 1, 2 & 2 & РЦНЯ \\
3 & 2, 3 & 2 & СНЭЯ \\
4 & 3, 4 & 1 & СОКЯ \\
5 & 4 & 1 & СОНС \\
\bottomrule
\end{tabular}
\end{table}

Четвёртый блок выбивается из общей картины: стык он пересекает, но изменяются у него
биты 13, 14 и~15 --- все три внутри третьей буквы, --- тогда как единственный его бит,
принадлежащий четвёртой букве, это бит~16, то есть уцелевший $d_4$. Следовательно, две
буквы портятся не всегда, когда блок расположен на стыке, а только когда граница букв
проходит внутри тех его битов, которые действительно изменились.

\subsection*{Вывод по заданию~3}
\addcontentsline{toc}{subsection}{Вывод по заданию~3}

Код Хэмминга $(7,4)$ построен, слово <<СОНЯ>> закодировано в 35 бит и восстановлено
после вмешательства. Обе матрицы согласованы: равенства $HG = 0$ и $RG = E$ проверены
прямым умножением.

Основной результат состоит в том, что код исправляет одну ошибку \emph{на блок}, а не
одну ошибку на сообщение. Обе серии опытов отличались только распределением ошибок по
блокам, а результаты получились противоположными. Четыре ошибки, разнесённые по
четырём разным блокам (серия B), исправлены все до одной, и слово восстановлено
полностью. Две ошибки, попавшие в один блок (серия A), код уже не исправил.
Следовательно, характеристика <<код исправляет одну ошибку>> относится к блоку из семи
бит, и на 35-битном сообщении из пяти блоков предел составляет пять исправляемых
ошибок --- но лишь при условии, что они распределены по одной на блок.

Это та же логика блочности, что и в задании~1: там искажение шифрограммы не выходило
за пределы своего блока, здесь ошибка исправляется или не исправляется независимо от
того, что происходит в соседних блоках. Различие в том, что в шифре Хилла размер блока
определялся выбранным размером ключа, а здесь он задан самим кодом.

Причина, по которой две ошибки не исправляются, видна из равенства~\eqref{eq:syn}.
Синдром есть сумма столбцов матрицы $H$ на позициях ошибок, и принимает он всего 8
значений, тогда как возможных векторов ошибок $2^7 = 128$. Однозначно распознаются
лишь 8 из них: нулевой и семь одиночных. Любой другой вектор ошибок даёт синдром,
совпадающий с синдромом какой-либо одиночной ошибки, и декодер выбирает именно её ---
как наиболее правдоподобный вариант из всех, дающих такой синдром. Разложение
$\vect{s} = H\vect{e}$ показывает и то, почему декодирование вообще возможно без знания
исходного сообщения: синдром зависит только от ошибки.

Отдельно следует назвать наиболее неблагоприятный сценарий --- три ошибки с нулевым
синдромом. В этом случае код не просто ошибается, но и не обнаруживает ошибку вовсе,
выдавая результат как достоверный. Происходит это тогда, когда вектор ошибок сам
оказывается кодовым словом. Минимальный вес ненулевого кодового слова в коде Хэмминга
$(7,4)$ равен трём (таких слов семь), откуда минимальное кодовое расстояние $d = 3$.
Общее правило: код исправляет $\lfloor (d-1)/2 \rfloor = 1$ ошибку и обнаруживает
$d - 1 = 2$ ошибки. Проведённый опыт с двумя ошибками --- как раз случай, когда ошибка
обнаружима (синдром ненулевой), но неисправима: декодер видит несоответствие, но
исправляет не ту позицию. Различать эти два режима код самостоятельно не умеет; чтобы
он только обнаруживал ошибки и не портил данные, его требуется принудительно переводить
в режим обнаружения.

Наконец, о контрольном вопросе, касающемся пространств. Столбцы матрицы $G$ порождают
линейную оболочку $\Imm G \subset \mathbb{F}_2^7$ --- множество всех 16 кодовых слов
размерности~4, поскольку $\rank G = 4$: строки 3, 5, 6 и~7 матрицы $G$ образуют
единичную матрицу $4\times4$, значит её столбцы линейно независимы. Ядро $\Ker H$ ---
множество векторов, проходящих все три проверки, и его размерность равна
$7 - \rank H$. Ранг $H$ равен трём, что также видно непосредственно из построения:
столбцы 1, 2 и~4 матрицы $H$ --- это $001$, $010$ и~$100$, вместе они составляют
единичную матрицу $3\times3$, поэтому три строки $H$ линейно независимы, а большего
числа строк в ней нет. Отсюда $\dim \Ker H = 7 - 3 = 4$.

Оба ранга вычисляются над полем $\mathbb{F}_2$, и это требует отдельной оговорки.
Здесь ответ над $\mathbb{Q}$ случайно совпал бы, но совпадать он не обязан: строки
$(1,1,0)$, $(0,1,1)$, $(1,0,1)$ над $\mathbb{Q}$ линейно независимы (их определитель
равен~2), тогда как над $\mathbb{F}_2$ их сумма равна нулю и ранг падает до двух. Ранг
есть свойство матрицы вместе с полем, а не матрицы самой по себе. Поскольку $HG = 0$,
справедливо включение $\Imm G \subseteq \Ker H$, а совпадение размерностей даёт
равенство
\begin{equation}
  \Imm G = \Ker H .
  \label{eq:imker}
\end{equation}
Матрицы $G$ и $H$ описывают одно и то же множество с двух сторон: $G$ его строит,
$H$ его распознаёт. Этим и объясняется, почему кодирование и проверка согласуются
автоматически и почему для подтверждения правильности матриц достаточно одной проверки
$HG = 0$.

\unnumsection{Заключение}

В работе матричное умножение применено к двум различным задачам, и обе решались не над
полем вещественных чисел, а над конечными структурами: кольцом вычетов
$\mathbb{Z}_{29}$ в заданиях~1 и~2 и полем $\mathbb{F}_2$ в задании~3. Это оказалось
не формальностью, а обстоятельством, определяющим всю технику работы. Обратимость
матрицы перестала сводиться к условию $\det \ne 0$ и превратилась в условие
$\gcd(\det K, n) = 1$; деление заменилось умножением на обратный элемент, который
приходится искать расширенным алгоритмом Евклида; стандартная функция
\texttt{numpy.linalg.inv} оказалась неприменимой. Выбор простого модуля $n = 29$ в
первых двух заданиях был сделан именно для того, чтобы условие обратимости упростилось
до $\det K \not\equiv 0$.

Все три задания оказались посвящены одному свойству --- блочности --- и тому, чем за
неё приходится платить. Матрица применяется к блоку целиком и независимо от соседних
блоков. В заданиях~1 и~2 это означает, что искажение не распространяется по сообщению,
но внутри своего блока портит, вообще говоря, всё. В задании~3 --- что ошибки
учитываются и исправляются поблочно, и четыре ошибки в четырёх блоках исправляются, а
две ошибки в одном блоке уже нет.

Сравнение двух методов показывает, что они устроены зеркально, и всё различие между
ними --- это различие между биекцией и вложением. Шифр Хилла биективен: матрица $K$
квадратная и обратимая, она отображает $\mathbb{Z}_{29}^z$ на себя целиком, длина
сообщения не меняется, информация не добавляется и не теряется. Код Хэмминга
инъективен, но не сюръективен: матрица $G$ прямоугольная и вкладывает $\mathbb{F}_2^4$
в $\mathbb{F}_2^7$ без малейшей потери --- именно это означает проверенное равенство
$RG = E$, исходный блок извлекается из кодового слова всегда и точно, --- однако образ
этого вложения занимает лишь 16 точек из 128, то есть на всё пространство $G$ не
отображает.

Из этого различия вытекают оба результата работы. У шифра Хилла образ --- всё
пространство, поэтому любая испорченная шифрограмма остаётся законной шифрограммой, и
обнаружить вмешательство невозможно в принципе: в задании~1 искажённые сообщения
расшифровались без каких-либо признаков порчи, просто в бессмысленный текст. Та же
биективность вместе с линейностью даёт и атаку из задания~2: ключ вычисляется из пары
<<открытый текст~--- шифрограмма>> по $z^2$ известным символам. У кода Хэмминга
наоборот: образ мал, и само попадание вне образа уже служит сигналом об ошибке.
Свободное место вокруг кодовых слов и есть тот ресурс, за счёт которого ошибку удаётся
не только заметить, но и локализовать; платой за него служит рост длины сообщения в
$7/4$ раза. Отсюда же виден и предел метода: свободного места ровно столько, сколько
необходимо на одну ошибку в блоке, а три ошибки переносят принятый вектор в соседнее
кодовое слово, то есть обратно внутрь образа, где сигнала нет и быть не может. Шифр
защищает от чтения и совершенно не защищает от искажений; код защищает от искажений и
ничего не скрывает.

Работа с синдромом показала, ради чего вообще строятся такие матрицы. Равенство
$\vect{s} = H(\vect{c} + \vect{e}) = H\vect{e}$ означает, что результат проверки не
зависит от передаваемых данных: декодер работает с ошибкой напрямую, не зная и не
восстанавливая исходное сообщение. Ограничение метода вытекает из того же равенства:
синдром принимает 8 значений при 128 возможных векторах ошибок, поэтому различимы
только нулевой вектор и семь одиночных ошибок. Все остальные случаи декодер вынужден
трактовать как ближайшую одиночную ошибку, а иногда --- когда вектор ошибок сам
оказывается кодовым словом --- не замечает ошибку вовсе.

\newpage
\unnumsection{Приложение А. Код}

\lstinputlisting[style=py,caption={Общий модуль \texttt{hill.py}},label={lst:hill}]{../code/hill.py}

\lstinputlisting[style=py,caption={Задание 1. Шифр Хилла},label={lst:task1}]{../code/task1.py}

\lstinputlisting[style=py,caption={Задание 2. Взлом шифра Хилла},label={lst:task2}]{../code/task2.py}

\lstinputlisting[style=py,caption={Задание 3. Код Хэмминга},label={lst:task3}]{../code/task3.py}

\end{document}
